\project{06}{The inertial unit as the joint's inner ear}{Motion sensing}{Rate and acceleration, FIFO, watermark, timestamp alignment with the control period}

\begin{keyfacts}
\item[Adds to the node] The first measurement of something physical in this volume: angular rate and acceleration, each sample carrying a timestamp the control loop can use
\item[Peripherals] The two-wire bus on the shield header, one data-ready pin routed to a capture channel from chapter 3, and one timer channel
\item[Depends on] Chapter 3 for the time base and the capture, chapter 4 for the pin budget, chapter 5 for the velocity this chapter cross-checks
\item[Real or modelled] \textbf{Real}, and the only fully real sensing in the volume. The shield is on the bench, it measures gravity and rotation, and nothing here is generated
\item[Difficulty] 3 of 5
\item[Effort] Three evenings of about three hours, one of them on the axis alignment
\item[Deliverable] A batched sensor read that never polls, a per-sample timestamp reconstructed from one hardware capture, the sensor's true sample rate measured in parts per million against the node's own clock, and a rotation from the sensor's axes to the joint's
\end{keyfacts}

\subsection*{Why this chapter}

An inertial unit on a joint is not a position sensor and pretending otherwise is
the fastest way to produce a demonstration that falls apart under questioning. A
single unit mounted on one link measures the angular rate of that link in space
and the acceleration it experiences, including gravity, in its own axes. It
cannot tell you the joint angle. Two units, one on each side of the joint, can
tell you a relative orientation, and even then only by integrating rate, which
drifts, or by leaning on gravity, which requires the joint to be still.

What it is genuinely good for on a joint node is a shorter and more useful list
than the marketing suggests. It sees vibration, which is what chatter and
resonance look like before they become visible. It sees a sudden contact, far
faster than a controller notices a following error. It gives an inclination when
the joint is at rest, which is an absolute reference that an incremental encoder
does not have. And it provides a second, independent opinion about whether the
link is moving at all, which is the beginning of the redundancy that chapter 18
needs.

The firmware subject of this chapter is none of those. It is **time**. The
sensor samples on its own oscillator, which is not the node's, batches samples
into a buffer, and raises a pin when the buffer reaches a watermark. By the time
the firmware reads that buffer, the oldest sample in it is tens of milliseconds
old, the newest is not, and neither carries a timestamp. Turning that into
samples the control loop can use is chapter 3's machinery applied to a real
sensor, and it is the part almost every project gets approximately right and
very few get exactly right.

\begin{note}{What this chapter does not repeat}
The sibling firmware volume already builds the driver layer for this shield in
full: the register-level drivers, the porting shim, the three bus topologies the
shield supports, the buffer and its watermark interrupt, and a ten minute run
with zero dropped samples counted on the board. That work is cited here and not
repeated. This chapter takes the driver as given and spends its pages on what a
joint node needs and that volume did not: the timestamp per sample, the measured
rate against the node's own clock, the axis alignment, and the cross-check
against the encoder from chapter 5.
\end{note}

\subsection*{Prior art and what to reuse}

\begin{table}[H]\centering\small
\begin{tabular}{p{34mm}p{46mm}p{38mm}p{20mm}}
\toprule
Source & What it gives & What it does not & Licence \\
\midrule
The silicon vendor's register-level driver collection for these sensors & Every sensor on both shields behind a two-function porting shim, permissively licensed throughout, so the register access is solved and reviewed & No board layer, no buffer policy, no timestamps, and no opinion about which axis is which & BSD-3-Clause \\
The sibling firmware volume, on this shield & The driver bring-up, the three bus topologies, the buffer with its watermark interrupt, and the ten minute zero-loss run. It is the chapter this one stands on & It is a sensor chapter, not a joint chapter: it never needs a sample to line up with a control period & this series \\
The sibling firmware volume, on filters and transforms & The filter this chapter's rate signal needs, already built and measured on this exact part, with the cycle cost stated & Nothing about sensor time & this series \\
The upstream operating system's description of this shield & A published machine-readable statement of the shield's bus topology and interrupt pins, which is the fastest way to settle which pin the data-ready line reaches & It describes the shield, not which of this board's timer channels that pin can reach, which is the question in step 2 & Apache-2.0 \\
The field oriented motor control project & The idiom for a sensor read inside a control loop, and a filter cascade that is tuned rather than theoretical & It reads position sensors, not inertial ones & MIT \\
\bottomrule
\end{tabular}
\caption{Prior art for chapter 6. The second row is the important one: this chapter is a layer on top of a chapter in another volume, and saying so is more honest than rebuilding the driver to make this volume look self-contained.}
\end{table}

What is left to write is the timestamp reconstruction, the rate measurement, the
axis alignment and the cross-check. No published source was found that
reconstructs per-sample timestamps for a batched inertial unit from a single
hardware capture on this part, so that construction is this volume's own.

\subsection*{What the node gains}

Before this chapter the node knows where it is, in counts, from a source it
generates itself. After it, the node measures something: the angular rate of the
link and the acceleration acting on it, in the joint's own axes, with each
sample carrying a timestamp on the same clock as everything else. It also gains
its first independent check, because the rate the inertial unit reports and the
velocity the encoder reports are now two opinions about the same motion.

\subsection*{Parts from the inventory}

\begin{table}[H]\centering\small
\begin{tabular}{p{40mm}p{66mm}p{40mm}}
\toprule
Part & Role & Interface \\
\midrule
NUCLEO-H7A3ZI-Q & The node & Micro USB to the host \\
One inertial and environmental shield & The sensor. \textbf{One shield at a time, never two}: the two on this bench share addresses and stacking them produces a bus that answers for the wrong part & Arduino header, two-wire bus \\
One jumper wire & The data-ready line to a timer capture channel, if the shield's interrupt pin does not already land on one & Header to header \\
A flat surface and a known right angle & The axis alignment in step 7 needs nothing more than a table and a book & None \\
\bottomrule
\end{tabular}
\caption{Inventory items used in chapter 6. The stacking rule comes from the bench notes and from the sibling volume, and it is repeated here because it is the kind of thing that costs an evening the first time.}
\end{table}

\subsection*{System architecture}

\diagram{j06_arch}{Where a sample comes from and what happens to it. Everything
solid is real. The single dashed block is the thing an inertial unit cannot give
a joint node, drawn so that nobody has to ask: one unit on one link does not
measure the joint angle, and the chapter says what it does measure instead.}

\subsection*{Peripheral configuration}

\begin{table}[H]\centering\small
\begin{tabular}{p{22mm}p{26mm}p{24mm}p{30mm}p{30mm}}
\toprule
Peripheral & Mode & Clock source & Pins and function & Interrupt and transfers \\
\midrule
I2C on the shield header & Fast mode, 400 kHz & Peripheral bus & Two pins, per the shield's topology & Transfers for the burst read, not the period handler \\
Timer capture channel & Input capture, rising edge & Timer clock, from chapter 3 & The sensor's data-ready line & Capture interrupt, low priority \\
The monotonic clock & From chapter 3, unchanged & As above & None & None \\
\bottomrule
\end{tabular}
\caption{Peripheral configuration for chapter 6. Whether the shield's data-ready pin lands on a pin this board can route to a timer capture channel is on chapter 4's confirm list, and step 2 gives the two fallbacks if it does not.}
\end{table}

\subsection*{Wiring}

\diagram{j06_wiring}{One shield, one rule and possibly one wire. The rule is the
one that costs an evening: the two shields on this bench must never be stacked
together. The wire is needed only if the data-ready pin does not already reach a
capture-capable pin, which is a board question rather than a shield question.}

\subsection*{Memory and timing budget}

\begin{table}[H]\centering\small
\begin{tabular}{p{44mm}p{28mm}p{28mm}p{30mm}}
\toprule
Quantity & Budget & Measured & Margin \\
\midrule
Flash, this chapter & 8 kB, excluding the vendor drivers & not measured & not measured \\
Static memory, this chapter & 2 kB, mostly the sample ring & not measured & not measured \\
Burst read duration & under 2 ms & not measured & not measured \\
Cost in the period handler & 0 cycles & by construction & n/a \\
Samples lost in one hour & 0 & not measured & not measured \\
Sensor rate against the node clock & report in ppm & not measured & not measured \\
Timestamp error, oldest sample & under 1 sample interval & not measured & not measured \\
\bottomrule
\end{tabular}
\caption{The budget for chapter 6. The fourth row is a design statement rather than a measurement: the period handler never touches the sensor, and a chapter that needed it to would be a chapter that had made a mistake.}
\end{table}

\subsection*{Firmware design (UML)}

\diagram{j06_uml}{One batch, from the sensor's own clock to the control loop's.
The capture on the data-ready edge is the only timestamp the hardware gives, and
the two self-messages in the middle are the whole of what this chapter adds:
every other sample's time is arithmetic on that one number, and the interval
used in that arithmetic is measured rather than taken from a datasheet.}

Three decisions.

\textbf{The period handler never talks to the sensor.} A two-wire transaction
takes hundreds of microseconds and can stall. The handler reads the most recent
aligned sample out of a ring buffer that somebody else filled, and that is all.

\textbf{The batch is read by a task, on the watermark interrupt.} The sensor
raises a pin when its buffer reaches a level; a low-priority task wakes, reads
the whole batch in one transfer, and stamps it. Reading one sample at a time on
every data-ready edge works and costs one bus transaction per sample, which at a
few hundred samples a second is a great deal of bus time for nothing.

\textbf{Neither signal corrects the other.} The rate from the inertial unit and
the velocity from the encoder are compared, the difference is reported, and
nothing is adjusted. The moment one of them is used to correct the other, the
node has one opinion again and the redundancy that chapter 18 wants has been
spent.

\subsection*{Data flow (ASCII)}

\begin{asciiart}
  sensor, on its own oscillator            node, on the clock from chapter 3
  +-----------------------------+          +-----------------------------------+
  | sample at the nominal rate  |          |                                   |
  |   |                         |          |                                   |
  |   v                         |          |                                   |
  | FIFO, N samples deep        |          |                                   |
  |   |                         |          |                                   |
  |   | watermark reached       |          |                                   |
  |   v                         |  pin     |                                   |
  | INT pin rises --------------+--------->| timer capture: ONE timestamp,     |
  |                             |          |   taken at the pin, in hardware   |
  |                             |          |     |                             |
  |                             |  burst   |     v                             |
  | FIFO read, N samples -------+--------->| task wakes, reads the whole batch |
  |                             |  read    |     |                             |
  +-----------------------------+          |     v                             |
                                           | reconstruct: newest = capture,    |
                                           |   older = capture - k * interval  |
                                           |     |                             |
                                           |     v                             |
                                           | rotate into the joint's axes      |
                                           |     |                             |
                                           |     v                             |
                                           | ring buffer <---- period handler  |
                                           |                    takes the most |
                                           |                    recent sample  |
                                           +-----------------------------------+
\end{asciiart}

\subsection*{Repository layout}

\begin{plaincode}
joint-node/
  board/  nucleo_h7a3zi_q.yaml            # + the shield pins, with evidence
  doc/    sensor-frames.md                # + this chapter: the axis convention
  src/
    bsp/  node/  control/  time/
    sense/
      encoder.c quadgen.c velocity.c
      imu.c      imu.h                    # + this chapter: batch read, stamps
      imu_rate.c imu_rate.h               # + this chapter: the measured rate
      frames.c   frames.h                 # + this chapter: sensor to joint
      crosscheck.c crosscheck.h           # + this chapter: two opinions
    act/  bus/  mw/  safety/  update/
  vendor/
    sensor-drivers/                       # + the vendor's register drivers,
                                          #   BSD-3-Clause, with their SPDX tags
  test/
    test_frames.c                         # + rotation arithmetic, on the host
    test_imu_stamps.c                     # + reconstruction, on the host
  host/  tools/  README.md
\end{plaincode}

\subsection*{Steps}

\begin{steps}
\item \textbf{Bring the shield up using the sibling volume, and stop there.} The
driver, the bus topology and the watermark interrupt are that volume's chapter,
and repeating them here would add pages and no knowledge. What this chapter
needs from it is one working call that returns a batch of raw samples and the
number of them.
\begin{ccode}
/* The interface this chapter needs from the driver layer. Everything behind it
   is the sibling volume's subject and the vendor's BSD-licensed drivers. */
int imu_read_batch(imu_raw_t *out, int max);   /* returns how many it read */
\end{ccode}

\item \textbf{Get the data-ready line onto a capture channel, or say why you
could not.} This is the whole reason the chapter can claim honest timestamps.
Three cases, in order of preference: the shield's interrupt pin already lands on
a pin that can be a timer capture input, in which case configure it; it does
not, but a jumper wire from the shield's pin to a capture-capable pin costs
nothing; or neither is possible, in which case the fallback is an external
interrupt that reads the monotonic clock, and the chapter reports the added
latency rather than hiding it.
\begin{plaincode}
preferred   shield INT pin -> timer capture channel   stamp taken in hardware
fallback 1  jumper wire to a capture-capable pin      same, plus one wire
fallback 2  external interrupt, read the clock        stamp late by the latency,
                                                      and the latency is reported
\end{plaincode}

\item \textbf{Reconstruct a timestamp for every sample from one capture.} The
hardware gives one timestamp per batch, at the moment the watermark was reached.
The newest sample in the batch is the one that caused it. Every older sample is
one nominal interval earlier than the one after it. That is the whole
reconstruction, and its error is bounded by how well the nominal interval
matches the sensor's true interval, which step 4 measures.
\begin{ccode}
/* imu.c: one hardware timestamp, N software ones. */
void imu_stamp_batch(imu_raw_t *s, int n, uint64_t capture_us, uint32_t interval_us)
{
    /* s[n-1] is the newest: it is the sample that reached the watermark. */
    for (int k = 0; k < n; k++) {
        int age = (n - 1) - k;                       /* 0 for the newest */
        s[k].t_us = capture_us - (uint64_t) age * interval_us;
    }
}
\end{ccode}

\item \textbf{Measure the sensor's true sample rate, and do not trust the
nominal one.} The sensor's rate comes from its own oscillator, whose tolerance
is in its datasheet and is far looser than a crystal's. Over a batch of a
thousand samples the difference between the nominal interval and the true one is
the whole timestamp error of the oldest sample. Measuring it costs nothing: the
node already has two hardware captures separated by a known number of samples.
\begin{ccode}
/* imu_rate.c: the interval the sensor actually achieves, in microseconds,
   measured against the node's own clock rather than assumed. */
void imu_rate_update(uint64_t capture_us, int samples_since_last)
{
    uint64_t dt = capture_us - g_last_capture_us;
    g_measured_interval_us = (float) dt / (float) samples_since_last;
    g_ppm = 1e6f * (g_measured_interval_us - g_nominal_interval_us)
                 / g_nominal_interval_us;
    g_last_capture_us = capture_us;
}
\end{ccode}
Report that figure in parts per million next to the oscillator comparison from
chapter 3, with the same caveat: it is a comparison of two clocks and it does
not say which is wrong. The difference is that here it does not matter, because
what the node wants is the sensor's rate expressed in the node's own time.

\item \textbf{Align the samples to the control period.} The control loop runs at
one rate and the sensor at another, and the two are not related by a whole
number. Two strategies, and the chapter builds both and states the error of
each.
\begin{ccode}
/* Nearest: take the sample whose timestamp is closest to the period boundary.
   Error: up to half a sample interval of age. Cost: nothing. */
const imu_sample_t *imu_nearest(uint64_t t_period_us);

/* Interpolated: linear between the two samples that straddle the boundary.
   Error: the curvature the straight line misses. Cost: a multiply and an add,
   and a signal that is no longer a measurement but an estimate. */
imu_sample_t imu_interpolate(uint64_t t_period_us);
\end{ccode}
The node uses the nearest sample and records its age with it, because a
controller that receives an estimate labelled as a measurement is a controller
whose error budget is wrong. The interpolated variant exists for the rig in
chapter 20, where a smooth signal is more useful than a raw one.

\item \textbf{Convert to physical units once, in one place.} Raw counts become
radians per second and metres per second squared exactly once, at the boundary
of the driver layer, and nothing above ever sees a raw count. The scale factor
comes from the configured full-scale range, which means the conversion has to
ask the driver what range is configured rather than carrying a constant.

\item \textbf{Find the rotation from the sensor's axes to the joint's.} The
shield is mounted however it fits, and its axes are not the joint's. With the
joint still, the unit measures gravity, and gravity points down. Three
orientations against a flat surface and a known right angle give three vectors,
and those give the rotation.
\begin{shellcode}
python host/frame_fit.py --flat readings_flat.csv \
                         --on-edge readings_edge.csv \
                         --on-end readings_end.csv -o doc/sensor-frames.md
# residual after fit: 0.4 degrees
# writing the rotation into board/nucleo_h7a3zi_q.yaml as sensor_to_joint
\end{shellcode}
Two honest notes on that number. A rotation fitted from gravity alone fixes two
of the three angles well and the third one poorly, because rotation about the
gravity vector does not change what gravity looks like: the third angle comes
from the joint's own motion, by turning the joint and seeing which sensor axis
reports the rate. And the result belongs in the board description from chapter
4, with its evidence marked as measured rather than documented.

\diagram{j06_data}{Three frames and one rotation. The shield is mounted however
it fits, so its axes are not the joint's, and three static orientations against
a flat surface give most of the rotation that relates them. Most, not all:
turning the shield about the gravity vector does not change what gravity looks
like, so the third angle carries no information in a static measurement and
comes from the joint's own motion instead.}

\item \textbf{Cross-check against the encoder, and report rather than correct.}
The joint turns; the encoder reports a velocity and the inertial unit reports an
angular rate about the joint axis. They should agree. Where they do not, the
difference is the interesting signal: a mechanical compliance between the
encoder and the link, a mounting that has shifted, or a sensor whose rate has
drifted.
\begin{ccode}
/* crosscheck.c: two opinions, one difference, and nothing corrected. */
void crosscheck_update(float encoder_rate, float imu_rate, uint64_t t_us)
{
    float d = encoder_rate - imu_rate;
    g_cross.last = d;
    if (fabsf(d) > g_cross.threshold) g_cross.exceedances++;
    /* No correction here, ever. Chapter 18 decides what an exceedance means. */
}
\end{ccode}

\item \textbf{Run for an hour and count what was lost.} The acceptance criterion
is zero dropped samples, counted on the board from the sensor's own buffer
status rather than inferred from the host.
\begin{shellcode}
python host/imu_soak.py --minutes 60
# batches: 7031   samples: 449,984   dropped: 0   overruns: 0
# measured interval: 4.8021 us/sample nominal 4.8000 -> +437 ppm
# timestamp error, oldest sample in batch: 0.31 sample intervals
\end{shellcode}
\end{steps}

\subsection*{Build, flash and debug}

\diagram{j06_timing}{Why the timestamps have to be reconstructed rather than
assumed. The sensor's batch arrives all at once and the samples in it are not
all from the same moment; the control period boundaries fall wherever they fall;
and the two rates drift apart because they come from different oscillators. The
lower band shows the error that the measured interval removes.}

\begin{shellcode}
cmake --build build -j && probe-rs run --chip STM32H7A3ZITx build/firmware.elf
python host/imu_soak.py --minutes 60
\end{shellcode}

\begin{note}{When every sample in the batch has the same timestamp}
The reconstruction in step 3 was skipped and the capture time was written into
every sample. It looks correct on a plot, it passes every smoke test, and it
gives the oldest sample in the batch an error of the whole batch duration, which
at a deep buffer is tens of milliseconds. A joint node that acts on a rate
signal timestamped tens of milliseconds late will oscillate, and the controller
will be suspected long before the timestamps are. A second signature, timestamps that drift
against the wall clock over an hour, is the nominal interval being used where
the measured one belongs.
\end{note}

\subsection*{Verification and acceptance criteria}

\begin{itemize}
\item One hour of continuous operation with zero dropped samples, counted from
the sensor's own buffer status on the board.
\item Every sample carries a timestamp on the node's monotonic clock, and the
oldest sample in a batch has a reconstructed timestamp within one sample
interval of the truth, verified by comparing a batch's first sample against a
capture taken deliberately at that moment.
\item The measured sample interval is reported in parts per million against the
nominal one, and the reconstruction uses the measured value.
\item The period handler performs no bus transaction, proven by a check that no
sensor call appears in its call graph.
\item The rotation from sensor axes to joint axes is in the board description,
marked as measured, with its residual stated and with the third angle's separate
determination described.
\item Turning the joint by hand produces an encoder velocity and an inertial
rate that agree within the stated threshold, and the difference is reported
every period rather than only when it is large.
\item A deliberately loosened mounting produces a visible increase in the
cross-check difference, which is the demonstration that the check is worth
having.
\end{itemize}

\subsection*{Variants}

\begin{table}[H]\centering\small
\begin{tabular}{p{24mm}p{26mm}p{44mm}p{22mm}p{20mm}}
\toprule
Axis & Variant & What changes & Cost & Built in full in \\
\midrule
Acquisition & Batch on a watermark & The baseline. One interrupt per batch and one bus transaction per batch & A buffer, and samples that are not fresh & Here \\
Acquisition & One sample per data-ready edge & Freshest possible data & One bus transaction per sample, and an interrupt per sample & Here, as the comparison \\
Acquisition & Continuous transfers & The bus engine fills a buffer with no processor involvement at all & Set-up complexity, and the coherency question the sibling volume settles & Chapter 20 \\
Timestamp & Hardware capture on the pin & The baseline, and the reason this chapter can claim what it claims & One timer channel & Here \\
Timestamp & Read the clock in the handler & Works anywhere, and is late by the interrupt latency & Accuracy, and the latency is reported & Here, as fallback 2 \\
Alignment & Nearest sample, age recorded & A measurement, honestly labelled, with a bounded age & Up to half an interval of age & Here \\
Alignment & Linear interpolation & A smooth signal at the period boundary & It is an estimate, and labelling it a measurement would be wrong & Here, for the rig \\
Processing & On the node & The baseline: rotate, convert, compare & Processor time in a task & Here \\
Processing & In the sensor & These parts carry programmable cores that can detect events without waking the node at all & A separate toolchain, and the sibling volume's subject & Nowhere. Cited, not rebuilt \\
Redundancy & Compare and report & The baseline. Two opinions stay two opinions & Nothing & Here \\
Redundancy & Fuse into one estimate & The usual answer, and it spends the redundancy: after fusion there is one opinion again & The independence chapter 18 wants & Nowhere, and the reason is in the design notes \\
\bottomrule
\end{tabular}
\caption{Variants for chapter 6. The last pair is the one worth arguing about: fusing an inertial rate with an encoder velocity gives a better signal and removes the ability to notice that they disagree, and a joint node wants the second property more than the first.}
\end{table}

\subsection*{Pitfalls}

\begin{itemize}
\item Stacking both shields. They share addresses, the bus answers for the wrong
part, and the symptom is plausible-looking data from a sensor that is not there.
\item Writing the capture time into every sample in the batch. The oldest sample
is then wrong by the whole batch duration.
\item Using the nominal sample interval. The sensor's oscillator is not a
crystal, and its tolerance is in the datasheet.
\item Reading the sensor from the period handler. A two-wire transaction takes
hundreds of microseconds and can stall, and the period handler has a deadline.
\item Treating the inertial unit as a joint angle sensor. It measures the link's
motion in space, and one unit on one link cannot give a joint angle at all.
\item Fitting the axis rotation from gravity alone and believing all three
angles. Rotation about the gravity vector is invisible to gravity.
\item Fusing the two velocity signals and then claiming redundancy. After fusion
there is one signal, and disagreement can no longer be detected.
\item Converting raw counts to physical units with a constant scale factor while
the full-scale range is configurable. The conversion asks the driver what range
is set.
\end{itemize}

\subsection*{Best practices applied}

\begin{itemize}
\item A chapter in another volume is cited and not rebuilt, and the interface
between them is one function.
\item The one hardware timestamp available is taken at the pin, and everything
else is arithmetic on it with a bounded and stated error.
\item A rate that could be assumed is measured instead, and the measurement is
reported in the same units and with the same caveat as chapter 3's.
\item A measurement stays labelled a measurement and an estimate stays labelled
an estimate, and the control loop is given the first.
\item Two independent signals are compared and neither is allowed to correct the
other, because the comparison is the point.
\item A calibration result goes into the board description with its evidence
marked as measured, alongside the rows whose evidence is a document.
\end{itemize}

\subsection*{Stretch goals}

\begin{itemize}
\item Put a second inertial unit on the other side of the joint and compute a
relative orientation. That is what a real installation does, and the drift of
the integrated difference over a minute is a figure worth publishing.
\item Detect contact from the acceleration signal and measure how many
milliseconds earlier it appears than the controller's following error does. That
comparison is the strongest single argument for having the sensor at all.
\item Run the spectrum from the sibling volume's transform chapter over the rate
signal while the joint is driven, and find the resonance. Chapter 7's plant
model can then be given that resonance, which makes it a better model.
\item Move the event detection into the sensor's own programmable core and
measure the energy saved, using the method from the sibling volume.
\end{itemize}

\subsection*{Roadmap and next steps}

Chapter 7 gives the node an output: the timer configuration for a motor bridge,
complementary outputs with dead time, a fault input, and a plant model standing
in for the motor that is not on this bench.

The published progression from here is the sibling volume's two sensor chapters
followed by its transform chapter, in that order, because the filter belongs
after the signal exists and not before. The vendor's driver collection is worth
reading rather than only linking against: it is small, register-level, and it
shows what each sensor's configuration actually costs. For the wider subject of
inertial sensing on a manipulator, the honest statement is that this volume does
not cover orientation estimation at all, and a reader who needs it should reach
for the literature on that subject rather than for this chapter.

\subsection*{Portfolio evidence}

\begin{itemize}
\item The one hour soak: batches, samples, zero dropped, and the measured
interval in parts per million.
\item The timestamp error figure for the oldest sample in a batch, which is the
number that distinguishes this implementation from the common one.
\item The cross-check trace with a deliberately loosened mounting, showing the
difference growing, which is a small and convincing demonstration that the
redundancy is real.
\item The frames document, with the rotation, its residual and the honest note
about which angle came from gravity and which from motion.
\end{itemize}

\subsection*{Sources}

Normative references:
\begin{itemize}
\item The datasheet for the inertial unit on the fitted shield, for the
oscillator tolerance, the full-scale ranges, the buffer depth and the data-ready
behaviour.
\item The shield's user manual, for the bus topology, the addresses and the
interrupt pin routing.
\item The board user manual UM2408, for which header pins the shield occupies
and which of them can reach a timer capture channel.
\end{itemize}

Reusable implementations:
\begin{itemize}
\item The silicon vendor's register-level driver collection for these sensors,
BSD-3-Clause.\newline
\url{https://github.com/STMicroelectronics/STMems_Standard_C_drivers}
\item The upstream operating system's description of this shield, Apache-2.0,
which settles the bus topology and the interrupt pins in machine-readable
form.\newline
\url{https://docs.zephyrproject.org/latest/boards/shields/index.html}
\item The field oriented motor control project, MIT, for the idiom of a sensor
read inside a control loop.\newline
\url{https://github.com/simplefoc/Arduino-FOC}
\end{itemize}
